\section{Which cells the canonical repair changes in an original copy}
\label{sec:repairspectrum}
The two terminal families give a precise description of the
intersection of each original copy with the canonical edit set
\[
 S_h=R_*\times\bigcup_{0\le p<m}C_h(p).
\]
Here $|S_h|=m^2$ is a count of physical cells. The classification
reveals that the same repair affects the two terminal forms with
different multiplicities.

\begin{proposition}[Repair-incidence spectrum]
\label{prop:repairspectrum}
Every type-I physical copy of $H_4$ in $A_h$ contains exactly one
cell from $S_h$. Every type-II copy contains exactly two cells from
$S_h$. The number of original copies meeting $S_h$ in one cell is
$2m^6$, and the number meeting it in two cells is
$\frac12(m-2)m^5$.
\end{proposition}
\begin{proof}
The unique row type belonging to $R_*$ in either terminal form is
the fourth row. In type I, the only chosen leaf column is the third
column $C_h(p)$; the two first columns are markers and the fourth
is the signed parent. Therefore precisely the fourth-row,
third-column physical cell belongs to $S_h$. In type II, both the
second and third columns are leaves, namely $C_h(q)$ and $C_h(p)$.
They are distinct since $q<p$. The first column is a marker and
the fourth is a signed parent. Thus exactly two cells of the copy
belong to $S_h$. The counting assertion follows from the exhaustive,
disjoint terminal classification and the weighted counts in
Section~\ref{sec:count}.
\end{proof}

\begin{corollary}[The canonical rectangle hits every old copy]
Every copy of $H_4$ present in $A_h$ has at least one entry changed
in the passage to $B_h$.
\end{corollary}
\begin{proof}
Both terminal families meet $S_h$ by
Proposition~\ref{prop:repairspectrum}, and there are no other
original copies.
\end{proof}

There are two distinct assertions in the proof of the sharp distance.
The corollary identifies a small set intersecting the originally
present copies; Lemma~\ref{lem:repair} proves that after all cells
of this set are changed simultaneously, the resulting physical
matrix is completely $H_4$-free. The latter requires the repaired
entry-table classification, because a cell edit can participate in
configurations absent from the original matrix. The arithmetic
repair certificate establishes freeness directly for the edited
matrix, rather than inferring it from the hitting corollary.

The spectrum also provides a quantitative accounting check.
If one sums, over original physical copies, the number of their
entries inside $S_h$, then the answer is
\begin{equation}\label{eq:repairincidence}
 2m^6+2\cdot\frac{m-2}{2}m^5
   =(3m-2)m^5.
\end{equation}
This number counts copy--edit-cell incidences with multiplicity;
it is deliberately different from both the number of copies
$\left(\frac52-\frac1m\right)m^6$ and the cardinality $m^2$
of the editing rectangle. Equation~\eqref{eq:repairincidence}
can be verified independently from the terminal parameterization.
